\begin{table}[h]
\caption{Transcriptome cohorts used. Cases are kidney biopsies in every cohort; controls are not. This is the design feature Section~\ref{sec:procurement} tests.}\label{tab:cohorts}
\footnotesize
\setlength{\tabcolsep}{4pt}
\begin{tabular}{@{}llllrrl@{}}
\toprule
Cohort & Role & Compart. & Platform & Cases & Controls & Control tissue \\
\midrule
GSE104948 & discovery & glom. & GPL22945/GPL24120 & 12 & 26 & donor, nephrec. \\
GSE142025 & discovery & cortex & Illumina HiSeq4000 & 27 & 9 & nephrec. \\
GSE294519 & discovery & tubulo. & GPL24676 & 23 & 13 & nephrec. \\
GSE30528 & discovery & glom. & GPL571 & 9 & 13 & nephrec. \\
GSE96804 & discovery & glom. & GPL17586 & 41 & 20 & nephrec. \\
GSE104954 & compartment & tubulo. & GPL22945/GPL24120 & 17 & 26 & donor, nephrec. \\
GSE30529 & compartment & tubulo. & GPL571 & 10 & 12 & nephrec. \\
\botrule
\end{tabular}
\footnotetext{Compartments: glom., glomerulus; tubulo., tubulointerstitium; cortex, whole renal cortex. Control tissue: nephrec., tumour nephrectomy; donor, living kidney donor. All cohorts are expressed on the same frozen 9,900-gene Entrez space. Cohorts marked \emph{compartment} are the second compartment of donors already represented by an ERCB or Woroniecka series; they enter the procurement and comparator analyses but not the cohort-composition sweep, which requires independent cohorts.}
\end{table}
